\chapter{Fundamentos Estatísticos e Inferenciais para Agentes Cognitivos}

\begin{quote}
\textit{``Todos os modelos estão errados, mas alguns são úteis.''} --- George E. P. Box
\end{quote}

\section{Probabilidade Bayesiana e Inferência Causal}

\subsection{Teorema de Bayes como Motor de Atualização de Crenças}

O \textbf{Teorema de Bayes} (Bayes, 1763; Laplace, 1774) é o fundamento matemático da atualização racional de crenças frente a evidências:

\begin{equation}
P(\theta | D) = \frac{P(D | \theta) \cdot P(\theta)}{P(D)}
\label{eq:bayes}
\end{equation}

No OpenCode, o \textbf{Confidence Ledger} (Capítulo 9) implementa uma atualização bayesiana simplificada via EMA (Exponential Moving Average):

\begin{equation}
C_{t+1}(a) = \alpha \cdot \text{score}_t + (1-\alpha) \cdot C_t(a)
\label{eq:ema_confidence}
\end{equation}

Onde $C_t(a)$ é a confiança no agente $a$ no tempo $t$, $\text{score}_t$ é o resultado da tarefa mais recente, e $\alpha \in (0,1)$ controla a taxa de adaptação.

\subsection{Inferência Causal e o \textit{do}-calculus}

Pearl (2009) \cite{pearl2009causality} revolucionou a inferência causal com o \textbf{do-calculus}:

\begin{equation}
P(Y | do(X=x)) \neq P(Y | X=x) \quad \text{(confundimento)}
\label{eq:do_calculus}
\end{equation}

O operador $do(X=x)$ representa uma \textbf{intervenção} — forçar $X$ ao valor $x$, quebrando as influências causais sobre $X$. Esta distinção é crucial para o \textbf{BehavioralGate}: não basta observar correlação entre ações do agente e resultados; é necessário verificar que a ação \textit{causou} o resultado.

\subsection{PRÁXIS 4.1 — Motor de Inferência Bayesiana}

\begin{lstlisting}[language=Python, caption={Atualização Bayesiana de confiança em agentes}]
import numpy as np
from scipy import stats

class BayesianTrustUpdater:
    """Atualização bayesiana da confiança usando Beta-Binomial."""
    
    def __init__(self, prior_alpha=2, prior_beta=2):
        self.alpha = prior_alpha
        self.beta = prior_beta
    
    def update(self, successes, failures):
        """Atualiza distribuição posterior Beta(alpha+successes, beta+failures)."""
        self.alpha += successes
        self.beta += failures
        return self.get_stats()
    
    def get_stats(self):
        posterior = stats.beta(self.alpha, self.beta)
        return {
            'mean': posterior.mean(),
            'std': posterior.std(),
            'credible_95': posterior.interval(0.95),
            'prob_better_than_05': 1 - posterior.cdf(0.5)
        }

# Uso no OpenCode: substitui EMA simples quando precisão bayesiana é necessária
updater = BayesianTrustUpdater(prior_alpha=2, prior_beta=2)
print("Prior (não-informativo):", updater.get_stats())
updater.update(successes=85, failures=15)
print("Após 100 tasks (85% sucesso):", updater.get_stats())
\end{lstlisting}

\section{Processos Estocásticos em Cadeias de Raciocínio}

\subsection{Cadeias de Markov e MDPs}

Uma \textbf{Cadeia de Markov} modela a transição entre estados onde o futuro depende apenas do presente:

\begin{equation}
P(S_{t+1} = s' | S_t = s, S_{t-1} = s_{t-1}, \ldots) = P(S_{t+1} = s' | S_t = s)
\label{eq:markov}
\end{equation}

Um \textbf{MDP} (Markov Decision Process) estende isso com ações e recompensas:

\begin{equation}
V^\pi(s) = \mathbb{E}\left[\sum_{t=0}^{\infty} \gamma^t R(s_t, a_t) \mid s_0 = s, \pi\right]
\label{eq:mdp}
\end{equation}

No OpenCode, a \textbf{Token Economy} (Capítulo 10) é essencialmente um MDP onde estados são níveis de stake dos agentes, ações são voluntariar-se para tarefas (via CFP), e recompensas são os fees recebidos (ou penalidades de slashing).

\subsection{MCMC e Amostragem para Exploração de Espaços de Agentes}

Métodos \textbf{MCMC} (Markov Chain Monte Carlo) — Metropolis-Hastings (1953), Hamiltonian Monte Carlo (Neal, 2011) — permitem amostrar de distribuições complexas. O OpenCode aplica estes princípios na \textbf{exploração do espaço de agentes}: quando o Attention Router precisa selecionar agentes para uma task, ele explora o espaço de agentes proporcionalmente ao Trust Score (análogo à amostragem da posterior).

\begin{thebibliography}{99}
\bibitem{pearl2009causality} Pearl, J. (2009). \textit{Causality: Models, Reasoning, and Inference} (2nd ed.). Cambridge University Press. \doiurl{10.1017/CBO9780511803161}
\bibitem{gelman2013bayesian} Gelman, A., et al. (2013). \textit{Bayesian Data Analysis} (3rd ed.). CRC Press.
\bibitem{wasserman2004all} Wasserman, L. (2004). \textit{All of Statistics}. Springer. \doiurl{10.1007/978-0-387-21736-9}
\bibitem{fisher1935design} Fisher, R. A. (1935). \textit{The Design of Experiments}. Oliver and Boyd.
\bibitem{kohavi2012online} Kohavi, R., et al. (2012). Online Controlled Experiments at Large Scale. \textit{KDD 2013}. \doiurl{10.1145/2487575.2488217}
\bibitem{neal2011mcmc} Neal, R. M. (2011). MCMC Using Hamiltonian Dynamics. In \textit{Handbook of Markov Chain Monte Carlo}. \doiurl{10.1201/b10905-6}
\bibitem{shannon1948mathematical} Shannon, C. E. (1948). A Mathematical Theory of Communication. \textit{Bell System Technical Journal}, 27(3), 379--423. \doiurl{10.1002/j.1538-7305.1948.tb01338.x}
\bibitem{cover2006elements} Cover, T. M., \& Thomas, J. A. (2006). \textit{Elements of Information Theory} (2nd ed.). Wiley. \doiurl{10.1002/047174882X}
\end{thebibliography}
